\documentclass[11pt]{article}
\usepackage{hyperref}
\usepackage{latexsym}
\usepackage{amsmath}
\usepackage{amssymb}
\usepackage{amsthm}
\usepackage{epsfig}
\usepackage{amsfonts}
\usepackage{braket}


\newcommand{\op}[1]{\operatorname{#1}}
\newcommand{\tinyspace}{\mspace{1mu}}

\newcommand{\abs}[1]{\lvert #1 \rvert}
\newcommand{\bigabs}[1]{\bigl\lvert #1 \bigr\rvert}
\newcommand{\Bigabs}[1]{\Bigl\lvert #1 \Bigr\rvert}
\newcommand{\biggabs}[1]{\biggl\lvert #1 \biggr\rvert}
\newcommand{\Biggabs}[1]{\Biggl\lvert #1 \Biggr\rvert}

\renewcommand{\natural}{\mathbb{N}}
\newcommand{\rev}{\textup{\tiny R}}

\newenvironment{mylist}[1]{\begin{list}{}{
	\setlength{\leftmargin}{#1}
	\setlength{\rightmargin}{0mm}
	\setlength{\labelsep}{2mm}
	\setlength{\labelwidth}{8mm}
	\setlength{\itemsep}{0mm}}}
	{\end{list}}
	
\newcommand{\handout}[5]{
  \noindent
  \begin{center}
  \framebox{
    \vbox{
      \hbox to 5.78in { {\bf CSE 309 Theory of Automata} \hfill #2 }
      \vspace{4mm}
      \hbox to 5.78in { {\Large \hfill #5  \hfill} }
      \vspace{2mm}
      \hbox to 5.78in { {\em #3 \hfill #4} }
    }
  }
  \end{center}
  \vspace*{4mm}
}

\newcommand{\lecture}[4]{\handout{#1}{#2}{#3}{#4}{#1}}

\newtheorem{theorem}{Theorem}
\newtheorem{corollary}[theorem]{Corollary}
\newtheorem{lemma}[theorem]{Lemma}
\newtheorem{observation}[theorem]{Observation}
\newtheorem{proposition}[theorem]{Proposition}
\newtheorem{definition}[theorem]{Definition}
\newtheorem{claim}[theorem]{Claim}
\newtheorem{fact}[theorem]{Fact}
\newtheorem{assumption}[theorem]{Assumption}

\topmargin 0pt
\advance \topmargin by -\headheight
\advance \topmargin by -\headsep
\textheight 8.9in
\oddsidemargin 0pt
\evensidemargin \oddsidemargin
\marginparwidth 0.5in
\textwidth 6.5in

\parindent 0in
\parskip 1.5ex
%\renewcommand{\baselinestretch}{1.25}

\begin{document}

\lecture{Problem Set 4}{\textit{Assigned: Friday, 17 Apr}}{Spring 2026}{Due: \textit{11:59 pm Monday, 4 May}}

\centerline{{\Large To facilitate grading and timely feedback}}
\centerline{{\Large please note that all submissions are through \href{http://gradescope.com}{Gradescope}.}}
\centerline{{\Large Solve each problem on a new page.}}
\centerline{{\Large Please clearly identify collaborator names and the resources used for each problem.}}

\begin{enumerate}

\item Assume that the formal system $F$ is sound, that the axioms and deduction rules
of $F$ can all be given explicitly, and that $F$ is capable of expressing any ordinary mathematical reasoning, for example, about Turing machines and proofs. In other words, $F$ is simply one of the usual systems to which G\"odel's Incompleteness Theorem applies.

Let $L = \{ x \, | \, x \textrm{ encodes a proof, in $F$, that } 1 + 1 = 3\}$.
\begin{enumerate}
\item Is $L$ the empty language?
\item Is your answer to $(a)$ provable in $F$?
\item Is $L$ decidable?
\item Is your answer to $(c)$ provable in $F$?
\item Describe a Turing machine $M$, such that the statement:
 
\emph{`$M$ runs in $O(2^n)$ time, where $n$ is the length of $M$'s input'} 

is independent of $F$~(i.e., neither provable nor disprovable in $F$).
\end{enumerate}

\item Show that for all functions $f : \mathbb{N} \mapsto \mathbb{N}$, we have \textsf{SPACE}$(f (n)) \subseteq \textsf{TIME}(2^{O(f(n))})$.

\item Let $G$ be an undirected graph with $n$ vertices. Then a \emph{Hamilton path} is a simple path in $G$ that visits each vertex once~(i.e., has $n$ vertices and $n-1$ edges), while a \emph{Hamilton cycle} is a simple cycle in $G$ that visits each vertex once (i.e., has $n$ vertices and $n$ edges). Let $HAMPATH$ and $HAMCYCLE$ be the problems of deciding whether $G$ has a Hamilton path and Hamilton cycle respectively, given $G$ as input.
\begin{enumerate}
\item Show that if $G$ has a Hamilton cycle, then $G$ also has a Hamilton path.
\item Give an example of a graph $G$ that has a Hamilton path but no Hamilton cycle.
\item Give a polynomial-time reduction from $HAMCYCLE$ to $HAMPATH$.
\item Give a polynomial-time reduction from $HAMPATH$ to $HAMCYCLE$.
\end{enumerate}

\item In the \emph{quadratic programming} ($QUADPROG$) problem, the input is a system of equalities and inequalities, each involving polynomials of degree at most $2$ (with integer coefficients) in $n$ real variables $x_1, \ldots , x_n$. The problem is to decide whether there exists an assignment to $x_1, \ldots , x_n$ that satisfies all the constraints simultaneously. As an example, the system
\begin{align*}
x_1 + x_2 &\leq 1 \\
x_1 &\geq 0 \\
x_2 &\geq 0 \\
4x_1x_2 &\geq 1
\end{align*}
can be satisfied by setting $x_1 = x_2 = 1/2$, but if we replaced the last inequality by $x_1x_2 \geq 1$, then the system would be unsatisfiable.
\begin{enumerate}
\item Show that $QUADPROG$ is \textsf{NP}-hard, by reduction from any problem that was already proved \textsf{NP}-hard in class.
\item What is a difficulty in showing that $QUADPROG \in \textsf{NP}$?
\end{enumerate}

\item We have $\mathsf{EXP} = \bigcup_k \mathsf{TIME}\left( 2^{n^k} \right)$, and that $\mathsf{NEXP} = \bigcup_k \mathsf{NTIME} \left( 2^{n^k} \right)$. Just as it is open problem
whether $\mathsf{P} = \mathsf{NP}$, it is also an open problem whether $\mathsf{EXP} = \mathsf{NEXP}$. Show that if $\mathsf{P} = \mathsf{NP}$, then $\mathsf{EXP} = \mathsf{NEXP}$ also. (\emph{Hint:} Given a language $L \in \mathsf{NEXP}$, can you come up with a modified language
$L' \in \mathsf{NP}$, such that $L \in \mathsf{EXP}$ if and only if $L' \in \mathsf{P}$?)

\end{enumerate}
\end{document}